\section{Метод Боголюбова--Митропольского}
\label{app:bm}

\subsection{Уравнение с квадратичной нелинейностью}
\label{sec:bm-quad}

Уравнение \eqref{eq:ode} с произвольной постоянной $c_0$ того же порядка, что и квадратичные члены, записывается так:
\begin{equation}
\label{eq:bm}
\ddot\delta+\omega^2\delta
=\varepsilon\Biggl(
c_0
+c_2\delta^2
+c_d(\dot\delta)^2
\Biggr),
\qquad
\omega=\omega_s,
\end{equation}
\begin{equation}
\label{eq:bm-coeff}
c_2=-\frac{\omega^2}{\eta_0}\Bigl(\eta_1-\frac{\eta_0}{\gamma_s^2}\Bigr),
\qquad
c_d=-\frac{1}{\gamma_s^2}.
\end{equation}
В главе~\ref{sec:mag-base} выполнено $c_0=0$. Если к $\tau$ добавляется постоянная $\Lambda$ одного порядка с $\delta^2$, в скобке \eqref{eq:ddelta-raw} появляется слагаемое $\Lambda$. Здесь $\kappa=K\cos\varphi_s$, и по \eqref{eq:ws} выполнено $\kappa\omega_\rf=-\omega^2/\eta_0$. Умножение на этот множитель даёт на правой части \eqref{eq:ode} постоянную $-\omega^2\Lambda/\eta_0$. Перекрёстный член $\Lambda\delta$ имеет третий порядок по амплитуде и отбрасывается, откуда
\begin{equation}
\label{eq:c0-L}
c_0=-\frac{\omega^2\Lambda}{\eta_0}.
\end{equation}
Параметр $\varepsilon$ отмечает порядок малости и в ответе полагается равным единице~\cite{Bogoliubov1961}.

Решение ищется в виде
\begin{equation}
\label{eq:ansatz}
\begin{aligned}
\delta&=a\cos\psi+\varepsilon u_1(a,\psi)+\varepsilon^2 u_2(a,\psi),\\
\dot a&=\varepsilon A_1+\varepsilon^2 A_2,
\qquad
\dot\psi&=\omega+\varepsilon B_1+\varepsilon^2 B_2.
\end{aligned}
\end{equation}
Функции $u_k$ периодичны по $\psi$ с периодом $2\pi$ и не содержат первой гармоники: коэффициенты при $\cos\psi$ и $\sin\psi$ равны нулю. Этим фиксирован смысл $a$ и $\psi$: вся первая гармоника собрана в $a\cos\psi$.

Подстановка в левую часть до порядка $\varepsilon$ считается явно. Ведущий член $a\cos\psi$ имеет производную
\begin{equation}
\label{eq:v0}
\frac{\mathrm{d}}{\mathrm{d}t}(a\cos\psi)
=-a\omega\sin\psi
+\varepsilon(A_1\cos\psi-a B_1\sin\psi)+O(\varepsilon^2).
\end{equation}
Здесь использовано $\dot a=\varepsilon A_1+O(\varepsilon^2)$ и $\dot\psi=\omega+\varepsilon B_1+O(\varepsilon^2)$. Производная поправки $\varepsilon u_1$ на том же порядке равна $\varepsilon\omega\,\partial_\psi u_1$, потому что $\dot a$ уже содержит $\varepsilon$, а $\dot\psi$ заменяется на $\omega$: следующий порядок есть $O(\varepsilon^2)$. Сумма
\begin{equation}
\label{eq:v1def}
\dot\delta
=-a\omega\sin\psi
+\varepsilon v_1+O(\varepsilon^2),
\qquad
v_1=A_1\cos\psi-a B_1\sin\psi+\omega\,\partial_\psi u_1.
\end{equation}
Вторая производная ведущего куска $-a\omega\sin\psi$ равна
\begin{equation}
\label{eq:d2v0}
\frac{\mathrm{d}}{\mathrm{d}t}(-a\omega\sin\psi)
=-a\omega^2\cos\psi
+\varepsilon(-A_1\omega\sin\psi-a\omega B_1\cos\psi)+O(\varepsilon^2).
\end{equation}
Вторая производная куска $\varepsilon v_1$ равна $\varepsilon\omega\,\partial_\psi v_1+O(\varepsilon^2)$. Раскрытие
\begin{equation}
\label{eq:dv1}
\partial_\psi v_1=-A_1\sin\psi-a B_1\cos\psi+\omega\,\partial_{\psi\psi}u_1
\end{equation}
даёт ещё один такой же вклад $-A_1\omega\sin\psi-a\omega B_1\cos\psi$. Два одинаковых вклада складываются, и коэффициент равен двум. Вместе с $\omega^2\cdot\varepsilon u_1$ член порядка $\varepsilon$ в $\ddot\delta+\omega^2\delta$ есть
\begin{equation}
\label{eq:lhs1}
\omega^2(\partial_{\psi\psi}u_1+u_1)-2\omega A_1\sin\psi-2a\omega B_1\cos\psi.
\end{equation}
На каждом следующем порядке $\varepsilon^k$ новые неизвестные $A_k$, $B_k$, $u_k$ входят в левую часть так же: $\dot a$ и $\dot\psi$ стоят в производной $a\cos\psi$ линейно, а $u_k$ входит в $\delta$ и во вторую производную линейно. Произведения уже найденных младших поправок переносятся в правую часть. Поэтому на порядке $\varepsilon^k$
\begin{equation}
\label{eq:ordk}
\omega^2(\partial_{\psi\psi}u_k+u_k)-2\omega A_k\sin\psi-2a\omega B_k\cos\psi=f_k,
\end{equation}
где $f_k$ --- коэффициент при $\varepsilon^k$ после этого переноса.

Правило гармоник следует из прямой подстановки: $\partial_{\psi\psi}\cos(m\psi)=-m^2\cos(m\psi)$, значит
\begin{equation}
\label{eq:harm-rule}
(\partial_{\psi\psi}+1)\cos(m\psi)=(1-m^2)\cos(m\psi).
\end{equation}
При $m\neq 1$ оператор обратим. Если правая часть уравнения $\omega^2(\partial_{\psi\psi}+1)u=F_m\cos(m\psi)$ содержит такую гармонику, коэффициент в $u$ равен $F_m/[\omega^2(1-m^2)]$. При $m=1$ выполнено $1-m^2=0$, оператор аннулирует $\cos\psi$ и $\sin\psi$, и эти члены правой части определяют $A_k$ и $B_k$, а в $u_k$ не входят.

На первом порядке правая часть \eqref{eq:bm} берётся на ведущем решении $\delta=a\cos\psi$, $\dot\delta=-a\omega\sin\psi$. Поправки порядка $\varepsilon$ внутри $\delta^2$ и $(\dot\delta)^2$ попадают в $\varepsilon^2$. Тождества $\cos^2\psi=(1+\cos 2\psi)/2$ и $\sin^2\psi=(1-\cos 2\psi)/2$ дают
\begin{equation}
\label{eq:f1}
\begin{aligned}
f_1
&=c_0+c_2 a^2\frac{1+\cos 2\psi}{2}
+c_d a^2\omega^2\frac{1-\cos 2\psi}{2}\\
&=c_0+\frac{a^2}{2}(c_2+c_d\omega^2)
+\frac{a^2}{2}(c_2-c_d\omega^2)\cos 2\psi.
\end{aligned}
\end{equation}
Гармоник $\cos\psi$ и $\sin\psi$ в \eqref{eq:f1} нет: их коэффициенты равны нулю. Левая часть \eqref{eq:lhs1} совпадает с $f_1$, поэтому коэффициенты при $\sin\psi$ и $\cos\psi$ слева тоже равны нулю. В \eqref{eq:lhs1} коэффициент при $\sin\psi$ равен $-2\omega A_1$, коэффициент при $\cos\psi$ равен $-2a\omega B_1$. При $\omega\neq 0$ и $a\neq 0$ оба коэффициента обращаются в нуль только при $A_1=0$ и $B_1=0$. Амплитуда на этом порядке постоянна, сдвига частоты нет.

Оставшиеся гармоники определяют $u_1=U_0+U_2\cos 2\psi$. Вторая производная постоянной равна нулю, так что $(\partial_{\psi\psi}+1)U_0=U_0$ и
\begin{equation}
\label{eq:U0}
\omega^2 U_0=c_0+\frac{a^2}{2}(c_2+c_d\omega^2).
\end{equation}
На гармонике $\cos 2\psi$ номер $m=2$. По правилу \eqref{eq:harm-rule} множитель $1-m^2$ равен $1-4=-3$, так что $\omega^2(\partial_{\psi\psi}+1)$ умножает коэффициент $U_2$ на $-3\omega^2$. Это произведение равно коэффициенту при $\cos 2\psi$ в \eqref{eq:f1}, поэтому
\begin{equation}
\label{eq:U2}
-3\omega^2 U_2=\frac{a^2}{2}(c_2-c_d\omega^2),
\qquad
U_2=-\frac{a^2(c_2-c_d\omega^2)}{6\omega^2}.
\end{equation}

Коэффициенты \eqref{eq:bm-coeff} подставляются явно. Сумма
\begin{equation}
\label{eq:c-plus}
c_2+c_d\omega^2
=-\frac{\omega^2}{\eta_0}\eta_1+\frac{\omega^2}{\gamma_s^2}-\frac{\omega^2}{\gamma_s^2}
=-\frac{\omega^2\eta_1}{\eta_0}.
\end{equation}
Два слагаемых $\omega^2/\gamma_s^2$ взаимно уничтожаются: первое возникает при раскрытии $(\omega^2/\eta_0)\cdot(\eta_0/\gamma_s^2)$ внутри $c_2$, второе равно $c_d\omega^2$. Деление \eqref{eq:U0} на $\omega^2$ даёт
\begin{equation}
\label{eq:U0-simp}
U_0=\frac{c_0}{\omega^2}-\frac{\eta_1 a^2}{2\eta_0}.
\end{equation}
Разность
\begin{equation}
\label{eq:c-minus}
c_2-c_d\omega^2
=-\frac{\omega^2\eta_1}{\eta_0}+\frac{2\omega^2}{\gamma_s^2},
\end{equation}
потому что $-c_d\omega^2=+\omega^2/\gamma_s^2$ добавляется к такому же члену из $(\omega^2/\eta_0)\cdot(\eta_0/\gamma_s^2)$. Подстановка в \eqref{eq:U2} даёт
\begin{equation}
\label{eq:U2-simp}
U_2=\frac{a^2}{6}\Bigl(\frac{\eta_1}{\eta_0}-\frac{2}{\gamma_s^2}\Bigr).
\end{equation}
При $\varepsilon=1$, с $A_1=B_1=0$ и с началом отсчёта, убирающим постоянную в $\psi$,
\begin{equation}
\label{eq:bm-sol}
\begin{aligned}
\delta
&=\frac{c_0}{\omega_s^2}-\frac{\eta_1 a^2}{2\eta_0}
+a\cos\psi
+\frac{a^2}{6}\Bigl(\frac{\eta_1}{\eta_0}-\frac{2}{\gamma_s^2}\Bigr)\cos 2\psi,\\
\psi&=\omega_s t.
\end{aligned}
\end{equation}
Сдвиг частоты появляется на следующем порядке метода из проекции на $\cos\psi$. Поправка этого порядка, пропорциональная $a^2$, в физической задаче требует члена $\delta^3$ в $\tau$ и здесь не выписывается.

\subsection{Уравнение с кубической нелинейностью}
\label{sec:bm-cubic}

Уравнение \eqref{eq:phi-ode} записывается так:
\begin{equation}
\label{eq:bm-c}
\ddot\varphi+\omega^2\varphi=\varepsilon c_3\varphi^3,
\qquad
\omega=\omega_s,
\qquad
c_3=\frac{\omega^2}{6}.
\end{equation}
Перенос члена $\omega^2\varphi^3/6$ из левой части \eqref{eq:phi-ode} в правую даёт $c_3\varphi^3$ при $\varepsilon=1$. Параметр $\varepsilon$ отмечает порядок малости и в ответе полагается равным единице, как в \S\ref{sec:bm-quad}.

Схема \eqref{eq:ansatz}, \eqref{eq:ordk} применяется к $\varphi$ вместо $\delta$. В этой подстановке $a$ есть коэффициент при $\cos\psi$ в фазе. Функция $u_1$ по условию схемы не содержит $\cos\psi$ и $\sin\psi$.

На первом порядке правая часть берётся на $\varphi=a\cos\psi$. Тождество $\cos 3\psi=4\cos^3\psi-3\cos\psi$ решается относительно куба:
\begin{equation}
\label{eq:cos3id}
\cos^3\psi=\frac{3\cos\psi+\cos 3\psi}{4}.
\end{equation}
Подстановка \eqref{eq:cos3id} даёт
\begin{equation}
\label{eq:f1c}
f_1
=c_3 a^3\cos^3\psi
=\frac{3c_3 a^3}{4}\cos\psi
+\frac{c_3 a^3}{4}\cos 3\psi.
\end{equation}
Коэффициент при $\sin\psi$ в \eqref{eq:f1c} равен нулю: разложение \eqref{eq:cos3id} его не содержит. В левой части \eqref{eq:ordk} коэффициент при $\sin\psi$ равен $-2\omega A_1$. При $\omega\neq 0$ равенство коэффициентов даёт $A_1=0$.

Коэффициент при $\cos\psi$ в \eqref{eq:f1c} равен $3c_3 a^3/4$. В левой части \eqref{eq:ordk} он равен $-2a\omega B_1$: по \eqref{eq:harm-rule} при $m=1$ оператор $\partial_{\psi\psi}+1$ аннулирует $\cos\psi$, а первая гармоника в $u_1$ запрещена анзацем \eqref{eq:ansatz}. При $a\neq 0$ и $\omega\neq 0$
\begin{equation}
\label{eq:B1c}
-2a\omega B_1=\frac{3c_3 a^3}{4},
\qquad
B_1=-\frac{3c_3 a^2}{8\omega}.
\end{equation}

На гармонике $\cos 3\psi$ номер $m=3$. Правило \eqref{eq:harm-rule} даёт $1-m^2=1-9=-8$, так что $\omega^2(\partial_{\psi\psi}+1)$ умножает коэффициент $U_3$ на $-8\omega^2$. Это произведение равно коэффициенту $c_3 a^3/4$ в \eqref{eq:f1c}, поэтому
\begin{equation}
\label{eq:U3c}
-8\omega^2 U_3=\frac{c_3 a^3}{4},
\qquad
U_3=-\frac{c_3 a^3}{32\omega^2},
\qquad
u_1=U_3\cos 3\psi.
\end{equation}
Подстановка $c_3=\omega^2/6$ сокращает $\omega^2$. В $B_1$ получается $3(\omega^2/6)a^2/(8\omega)=\omega a^2/16$, знак минус сохраняется. В $U_3$ получается $(\omega^2/6)a^3/(32\omega^2)=a^3/192$, знак минус сохраняется:
\begin{equation}
\label{eq:B1U3}
B_1=-\frac{\omega a^2}{16},
\qquad
U_3=-\frac{a^3}{192}.
\end{equation}
При $\varepsilon=1$ частота есть $\dot\psi=\omega+B_1$, и
\begin{equation}
\label{eq:phi-a}
\varphi
=a\cos\psi-\frac{a^3}{192}\cos 3\psi,
\qquad
\Omega=\omega\Bigl(1-\frac{a^2}{16}\Bigr).
\end{equation}
Здесь $a$ --- коэффициент первой гармоники, не полуразмах.

В главе~\ref{sec:mag-phase} амплитуда задана полуразмахом; следующие равенства только переводят $a$ в $\hat\varphi$. При $\psi=0$ выполнено $\cos\psi=\cos 3\psi=1$, поэтому значение \eqref{eq:phi-a} равно $a-a^3/192$. При $\psi=\pi$ оно противоположно. Производная по $\psi$ равна $-a\sin\psi+(a^3/64)\sin 3\psi$: множитель $3/192$ равен $1/64$. Вынос $\sin\psi$ оставляет множитель $1-(a^2/64)(\sin 3\psi/\sin\psi)$. Отношение $\sin 3\psi/\sin\psi=3-4\sin^2\psi$ ограничено, и при $a^2\ll 1$ этот множитель в нуль не обращается. Нули производной только при $\sin\psi=0$. Полуразмах равен значению при $\psi=0$:
\begin{equation}
\label{eq:phihat-a}
\hat\varphi=a\Bigl(1-\frac{a^2}{192}\Bigr).
\end{equation}
Обращение с отбрасыванием $O(a^5)$ есть $a=\hat\varphi+\hat\varphi^3/192$. Подстановка этой формулы обратно в $a-a^3/192$ возвращает $\hat\varphi$, потому что $a^3=\hat\varphi^3+O(\hat\varphi^5)$. Подстановка того же $a$ в \eqref{eq:phi-a} даёт
\begin{equation}
\label{eq:phi-hat-cos}
\varphi
=\hat\varphi\cos\psi
+\frac{\hat\varphi^3}{192}\bigl(\cos\psi-\cos 3\psi\bigr),
\qquad
\frac{\Omega}{\omega_s}=1-\frac{\hat\varphi^2}{16}.
\end{equation}
В частоте $a^2=\hat\varphi^2+2\hat\varphi\cdot(\hat\varphi^3/192)+O(\hat\varphi^6)=\hat\varphi^2+O(\hat\varphi^4)$, поэтому замена $a^2$ на $\hat\varphi^2$ не меняет член $\hat\varphi^2/16$. Остаток в $\varphi$ есть $O(\hat\varphi^5)$.

